% Tabla 14: Ancho relativo del intervalo según la profundidad y el conocimiento disponible.
% Etiqueta original: tab:intervalo.
% Fuente: capitulos/4-resultados.tex, lineas 195-212.
% Fragmento extraido automaticamente; requiere el preambulo de main.tex.

\begin{table}[htbp]
\centering\footnotesize
\caption{Ancho relativo del intervalo $\IL$ (\% de la LT de calibre) según la profundidad de la réplica y el conocimiento disponible sobre $z$: cálculo numérico con C sobre las 72 capturas y, entre paréntesis, predicción geométrica de primer orden. La última columna indica la cobertura empírica.}
\label{tab:intervalo}
\begin{tabular}{@{}lccccc@{}}
\toprule
Conocimiento de $z$ & 5~cm & 11~cm & 17~cm & 23~cm & Cobertura \\
\midrule
Desconocida (0--30~cm) & 81,8 (83,3) & 71,0 (71,4) & 62,4 (62,5) & 55,6 (55,5) & 1,00 \\
Conocida con $\pm 3$~cm & 16,4 (16,7) & 14,2 (14,3) & 12,5 (12,5) & 11,1 (11,1) & 1,00 \\
Conocida con $\pm 1$~cm & 5,5 (5,6) & 4,7 (4,8) & 4,2 (4,2) & 3,7 (3,7) & 0,81 \\
\midrule
MAPE de \Aprime{} ($z$ conocida) & \multicolumn{4}{c}{2,3} & \\
MAPE de B ($z$ conocida) & \multicolumn{4}{c}{0,7} & \\
MAPE de C ($z$ conocida) & \multicolumn{4}{c}{0,9} & \\
\bottomrule
\end{tabular}
\end{table}
